\documentclass[11pt]{article}
\usepackage[margin=1in]{geometry}
\usepackage{booktabs}

\title{Project 5 Brief Performance Report}
\author{}
\date{}

\begin{document}
\maketitle

\section*{Implementation Summary}
\textbf{Task 1 (Vector Add):} Implemented CUDA vector addition with static device arrays, CUDA event timing, and bandwidth calculations (theoretical and measured). Correctness was verified with host-side checking (0 errors).\\
\textbf{Task 2 (Ray Tracing):} Completed the ray--sphere kernel using per-pixel intersection tests across all spheres, selected visible sphere by depth rule from the project formulation, and computed color via intersection ratio. Added lightweight CPU vs GPU verification: full compare for small sphere counts and sampled checks for larger counts.

\section*{Performance Comparison: RTX 4090 vs A100}

\subsection*{Task 1}
\begin{table}[h!]
\centering
\begin{tabular}{lcc}
\toprule
Metric & RTX 4090 (Local) & NVIDIA A100 80GB PCIe \\
\midrule
Kernel time (ms) & 0.019 & 0.036 \\
Measured BW (GB/s) & 2409.28 & 1306.73 \\
Theoretical BW (GB/s) & 1008.10 & 1935.36 \\
Measured / Theoretical & 238.99\% & 67.52\% \\
\bottomrule
\end{tabular}
\end{table}

\noindent The A100 result is plausible and stable. The 4090 value exceeds theoretical DRAM bandwidth because this tiny kernel can be dominated by timing granularity/overheads and cache effects at very short durations.

\subsection*{Task 2}
\begin{table}[h!]
\centering
\begin{tabular}{rcc}
\toprule
Spheres & RTX 4090 time (ms) & A100 time (ms) \\
\midrule
16   & 0.077 & 0.104 \\
32   & 0.105 & 0.148 \\
64   & 0.187 & 0.268 \\
128  & 0.345 & 0.486 \\
256  & 0.664 & 0.942 \\
512  & 1.305 & 1.845 \\
1024 & 2.573 & 3.654 \\
2048 & 5.124 & 7.272 \\
\bottomrule
\end{tabular}
\end{table}

\noindent Both GPUs show the expected near-linear scaling with sphere count. In this environment, the 4090 is consistently faster for this workload, while the A100 still produces correct outputs and stable timing trends.

\section*{Correctness Notes}
Task 1 reports 0 verification errors. Task 2 CPU/GPU checks show exact agreement for small cases except for rare \(\pm1\) channel differences at a few pixels (normal floating-point rounding behavior across architectures).

\end{document}
